\section{Síntesis y ampliación}

La aportación más importante de esta clase no es clasificar DPO, PPO o QLoRA, sino ofrecer una historia interpretable de los sistemas abiertos de alineación: cada salto de los modelos dependió a la vez del base model, de la distribución de los datos, del acceso al entrenamiento, de la retroalimentación de la evaluación y de los artifacts publicados. Solo al entender esta cadena se puede juzgar qué cuello de botella resuelve realmente un método nuevo.

\subsection{Catorce conclusiones centrales}

Esta sección condensa el curso en conclusiones reutilizables; cada una conserva sus condiciones y su nivel de evidencia.

\begin{enumerate}
  \item RLHF es importante para el ChatGPT-like behavior, pero no es condición suficiente para un producto completo.
  \item El base model es la base de las capacidades; la alignment principalmente reordena y moldea la distribución de las interacciones.
  \item IFT, SFT, RLHF y DPO deben distinguirse por sus datos y su función de pérdida.
  \item Alpaca demostró que la synthetic instruction data permite adaptar la interfaz con rapidez.
  \item Vicuna/ShareGPT demostró que la prompt distribution de usuarios reales es sumamente valiosa.
  \item OpenAssistant muestra que los datos humanos abiertos tienen un valor de reutilización a largo plazo.
  \item QLoRA amplió el número de participantes, pero no resolvió automáticamente la estabilidad de RLHF.
  \item El debate entre safety y uncensored es, en esencia, una elección sobre el límite de las negativas y el control del usuario.
  \item Arena, AlpacaEval, MT-Bench y Leaderboard sirven a plazos de retroalimentación distintos.
  \item LLM-as-a-judge es rápido, pero tiene sesgos de length, position, self-preference y contamination.
  \item El objetivo de RLHF combina reward y KL constraint; PPO es solo una de las vías de implementación.
  \item DPO es sencillo y eficaz, pero difiere de PPO en la dinámica de entrenamiento y en el aprovechamiento de los datos.
  \item Zephyr/Tulu demostraron que DPO puede llevarse a la práctica; SteerLM/Starling demostraron que RL sigue siendo competitivo.
  \item En 2024, el cuello de botella persistente de la alignment abierta es la preference data diversa, legal y publicable.
\end{enumerate}

\subsection{Una tabla que articula el sistema abierto de alineación}

Para no quedarse solo con nombres de modelos, a continuación se resume, por componente del sistema, el problema que resolvió cada etapa y el riesgo residual.

\begin{table}[H]
\centering
\small
\begin{tabularx}{\textwidth}{L{0.16\textwidth} L{0.22\textwidth} X}
\toprule
Componente & Artifact representativo & Problemas aún sin resolver \\
\midrule
Base model & LLaMA/Llama 2/Llama 3 & Capacidades, licencia, provenance de los datos y costo de entrenamiento \\
IFT data & Alpaca, ShareGPT, UltraChat & Intent real, diversidad, privacidad y sesgo de plantilla \\
Human data & OpenAssistant/Arena logs & Costo, consistencia, consent y cobertura \\
Efficient tuning & LoRA/QLoRA & Conservación completa de las capacidades, RL stability y fusión para inference \\
Evaluation & Arena/AlpacaEval/MT-Bench/Leaderboard & Judge bias, contaminación, latencia y gaming \\
Preference optimization & PPO/DPO/SteerLM & Eficiencia de datos, estabilidad, reward hacking e interpretabilidad \\
Ecosystem & AI2, HF, LMSYS, community & Lineage, reproducibilidad, mantenimiento a largo plazo y gobernanza \\
\bottomrule
\end{tabularx}
\caption{Componentes de la alineación abierta, artifacts representativos y riesgos residuales.}
\end{table}

\subsection{Release checklist para la práctica}

La historia del curso puede convertirse en una checklist previa a la publicación de un modelo. El objetivo es que «usamos DPO» deje de sustituir a la evidencia completa.

\begin{lstlisting}[caption={Open aligned model release checklist}]
record_base_model_and_license()
publish_chat_template_and_training_objective()
document_instruction_and_preference_data_provenance()
report_compute_memory_and_random_seeds()
run_capability_safety_and_human_preference_evals()
disclose_judge_models_baselines_and_known_biases()
release_failure_examples_and_model_card()
\end{lstlisting}

\begin{importantbox}{Cinco filtros para juzgar un nuevo método de alignment}
\begin{enumerate}
  \item ¿Supera a un baseline fuerte con base/data/eval compartidos?
  \item ¿Produce un modelo competitivo descargable y reproducible, y no solo curvas de entrenamiento?
  \item ¿Conserva las base capabilities y reduce la failure rate real?
  \item ¿Se sostiene con varios judges y con preferencias humanas, y no en un único leaderboard?
  \item ¿La ganancia justifica los datos, el rollout, la memoria de GPU y la complejidad de ingeniería adicionales?
\end{enumerate}
\end{importantbox}

\subsection{Preguntas de autoevaluación}

Las siguientes preguntas sirven para comprobar si se domina la causalidad del sistema, y no solo la cronología de los releases.

\begin{enumerate}
  \item ¿Por qué RLHF es importante para ChatGPT y, aun así, no es condición suficiente?
  \item ¿Cómo se distinguen IFT y SFT en objetivo, datos e interfaz?
  \item ¿Por qué ShareGPT puede ser más valioso que una mayor cantidad de synthetic prompts?
  \item ¿Cuál es la diferencia esencial entre los weight differences y una LoRA update?
  \item ¿Qué memoria ahorra QLoRA y qué costos de RLHF no ahorra?
  \item ¿Para qué decisiones sirve cada uno de Arena, AlpacaEval, MT-Bench y Leaderboard?
  \item ¿Qué evita la KL penalty en el RLHF objective?
  \item ¿Qué representa el log-ratio policy/reference en la DPO loss?
  \item ¿Por qué Zephyr/Tulu no bastan por sí solos para demostrar que DPO siempre supera a PPO?
  \item ¿Por qué la preference data sintética puede ampliar la escala y, a la vez, estrechar la distribución?
\end{enumerate}

\subsection{Lecturas adicionales}

\begin{itemize}
  \item Ouyang et al., \emph{Training language models to follow instructions with human feedback}: InstructGPT y el pipeline de RLHF.
  \item Taori et al., \emph{Stanford Alpaca}; Wang et al., \emph{Self-Instruct}: la synthetic instruction data temprana.
  \item Dettmers et al., \emph{QLoRA}: base model cuantizado y fine-tuning de bajo rango.
  \item Rafailov et al., \emph{Direct Preference Optimization}: derivación y experimentos de DPO.
  \item Bai et al., \emph{Constitutional AI}; Köpf et al., \emph{OpenAssistant Conversations}: preferencias y datos humanos abiertos.
  \item Zheng et al., \emph{Chatbot Arena}: pairwise evaluation humana y el ecosistema de modelos abiertos.
\end{itemize}

\end{document}
